% Seminar 3 - Page 5. Fragment for
% % Seminar 3 - Page 5. Fragment for
% % Seminar 3 - Page 5. Fragment for
% % Seminar 3 - Page 5. Fragment for
% \input{seminars/pages/seminar3-pagina5-en}; no preamble or page break.
% EDITORIAL NOTES:
% - The opening lines continue the universal property of the quotient set
%   from Page 4. The word attached to the second criterion is overwritten:
%   it may read "inj." or "bij.". The injectivity reading is provisional;
%   if "bij." was intended, surjectivity is also required.
% - The example defining rational numbers writes Z x Z without excluding
%   zero denominators. On all of Z x Z, cross-multiplication is NOT an
%   equivalence relation. The necessary restriction is not silently added.
% - The Cauchy-sequence construction is read as sequences with terms in Q;
%   it is written compactly in the original manuscript.
% - Example (i) includes a small sketch of the quotient projection, here
%   rendered as a commutative diagram.
% - The right-margin scribbles and small shaded circle are not part of an
%   identifiable calculation and are omitted from the main transcription.
% - Example (iv) appears to identify R/~ with R; this is only an abstract
%   (non-canonical) identification, not a specified quotient map.

\subsection*{Seminar 3 - Page 5}
\begingroup
\setlength{\abovedisplayskip}{5pt}
\setlength{\belowdisplayskip}{5pt}
\setlength{\abovedisplayshortskip}{3pt}
\setlength{\belowdisplayshortskip}{3pt}

\noindent\textbf{Universal property (continued).}
For the map $\overline f$ of the preceding page,
\[
\begin{aligned}
 \overline f\text{ is surjective}
   &\iff f\text{ is surjective},\\
 \overline f\text{ is injective}
   &\iff {\sim}\;=R_f.
\end{aligned}
\]
Moreover, for a system of representatives $(x_i)_{i\in I}$,
\[
 M=\bigcup_{i\in I}\widehat{x_i}.
\]

\medskip
\noindent\underline{Examples of quotient sets}

\noindent\textbf{(a) Rational numbers.}
Consider $\mathbb Z\times\mathbb Z$ and the relation
\[
 (m,n)\sim(m',n')
 \overset{\mathrm{def}}{\iff}mn'=nm'.
\]
The notes identify the quotient as
\[
 (\mathbb Z\times\mathbb Z)/{\sim}
   \overset{\mathrm{def}}{=}\mathbb Q,
 \qquad
 r\in\mathbb Q,\quad r=\widehat{(m,n)}.
\]
% Source caveat: n,n' != 0 is required for the intended construction.

\noindent\textbf{(b) Real numbers from Cauchy sequences.}
Let
\[
 \mathcal C=\bigl\{(x_n)_{n\geq1}\text{ with }x_n\in\mathbb Q
       \mid (x_n)_{n\geq1}\text{ is Cauchy}\bigr\}.
\]
Define
\[
 (x_n)_n\sim(y_n)_n
 \overset{\mathrm{def}}{\iff}
 \lim_{n\to\infty}(x_n-y_n)=0.
\]
This is an equivalence relation; with the manuscript's notation,
\[
 \mathcal C/{\sim}\ \overset{\mathrm{not.}}{=}\ \mathbb R,
 \qquad
 r=\widehat{(x_n)_n}\in\mathbb R.
\]

\medskip
\noindent\textbf{(i)} On $\mathbb R$, put
\[
 x\sim y\iff x-y\in\mathbb Z.
\]
This is an equivalence relation and\enspace\(\mathbb R/{\sim}\cong[0,1)\).
The quotient sketch is
\[
\begin{tikzcd}[column sep=large,row sep=small]
 \mathbb R\arrow[r,"p_{\sim}"]\arrow[d,"{x\mapsto x-\lfloor x\rfloor}"']
    & \mathbb R/{\sim} \\
 {[0,1)}\arrow[ur,dashed,"\cong"'] &
\end{tikzcd}
\]
% The floor-map label explicates the diagram's unlabeled vertical arrow.

\noindent\textbf{(ii)} On $\mathbb C$, define
\[
 z_1\sim z_2\overset{\mathrm{def}}{\iff}|z_1|=|z_2|.
\]
This is an equivalence relation on $\mathbb C$, with\enspace\(\mathbb C/{\sim}\cong[0,+\infty)\).

\noindent\textbf{(iii)} On $\mathbb C^*$, define
\[
 z_1\sim z_2\overset{\mathrm{def}}{\iff}
 \arg(z_1)=\arg(z_2).
\]
This is an equivalence relation on $\mathbb C^*$ and
\[
 \mathbb C^*/{\sim}\;\cong\;S^1,
 \qquad S^1=\{z\in\mathbb C\mid |z|=1\}.
\]
% Equality of arguments is understood with a consistent branch/modulo 2pi.

\noindent\textbf{(iv)} On $\mathbb R$, put
\[
 x\sim y\iff x-y\in\mathbb Q.
\]
The manuscript records $\mathbb R/{\sim}\cong\mathbb R$.

\endgroup
; no preamble or page break.
% EDITORIAL NOTES:
% - The opening lines continue the universal property of the quotient set
%   from Page 4. The word attached to the second criterion is overwritten:
%   it may read "inj." or "bij.". The injectivity reading is provisional;
%   if "bij." was intended, surjectivity is also required.
% - The example defining rational numbers writes Z x Z without excluding
%   zero denominators. On all of Z x Z, cross-multiplication is NOT an
%   equivalence relation. The necessary restriction is not silently added.
% - The Cauchy-sequence construction is read as sequences with terms in Q;
%   it is written compactly in the original manuscript.
% - Example (i) includes a small sketch of the quotient projection, here
%   rendered as a commutative diagram.
% - The right-margin scribbles and small shaded circle are not part of an
%   identifiable calculation and are omitted from the main transcription.
% - Example (iv) appears to identify R/~ with R; this is only an abstract
%   (non-canonical) identification, not a specified quotient map.

\subsection*{Seminar 3 - Page 5}
\begingroup
\setlength{\abovedisplayskip}{5pt}
\setlength{\belowdisplayskip}{5pt}
\setlength{\abovedisplayshortskip}{3pt}
\setlength{\belowdisplayshortskip}{3pt}

\noindent\textbf{Universal property (continued).}
For the map $\overline f$ of the preceding page,
\[
\begin{aligned}
 \overline f\text{ is surjective}
   &\iff f\text{ is surjective},\\
 \overline f\text{ is injective}
   &\iff {\sim}\;=R_f.
\end{aligned}
\]
Moreover, for a system of representatives $(x_i)_{i\in I}$,
\[
 M=\bigcup_{i\in I}\widehat{x_i}.
\]

\medskip
\noindent\underline{Examples of quotient sets}

\noindent\textbf{(a) Rational numbers.}
Consider $\mathbb Z\times\mathbb Z$ and the relation
\[
 (m,n)\sim(m',n')
 \overset{\mathrm{def}}{\iff}mn'=nm'.
\]
The notes identify the quotient as
\[
 (\mathbb Z\times\mathbb Z)/{\sim}
   \overset{\mathrm{def}}{=}\mathbb Q,
 \qquad
 r\in\mathbb Q,\quad r=\widehat{(m,n)}.
\]
% Source caveat: n,n' != 0 is required for the intended construction.

\noindent\textbf{(b) Real numbers from Cauchy sequences.}
Let
\[
 \mathcal C=\bigl\{(x_n)_{n\geq1}\text{ with }x_n\in\mathbb Q
       \mid (x_n)_{n\geq1}\text{ is Cauchy}\bigr\}.
\]
Define
\[
 (x_n)_n\sim(y_n)_n
 \overset{\mathrm{def}}{\iff}
 \lim_{n\to\infty}(x_n-y_n)=0.
\]
This is an equivalence relation; with the manuscript's notation,
\[
 \mathcal C/{\sim}\ \overset{\mathrm{not.}}{=}\ \mathbb R,
 \qquad
 r=\widehat{(x_n)_n}\in\mathbb R.
\]

\medskip
\noindent\textbf{(i)} On $\mathbb R$, put
\[
 x\sim y\iff x-y\in\mathbb Z.
\]
This is an equivalence relation and\enspace\(\mathbb R/{\sim}\cong[0,1)\).
The quotient sketch is
\[
\begin{tikzcd}[column sep=large,row sep=small]
 \mathbb R\arrow[r,"p_{\sim}"]\arrow[d,"{x\mapsto x-\lfloor x\rfloor}"']
    & \mathbb R/{\sim} \\
 {[0,1)}\arrow[ur,dashed,"\cong"'] &
\end{tikzcd}
\]
% The floor-map label explicates the diagram's unlabeled vertical arrow.

\noindent\textbf{(ii)} On $\mathbb C$, define
\[
 z_1\sim z_2\overset{\mathrm{def}}{\iff}|z_1|=|z_2|.
\]
This is an equivalence relation on $\mathbb C$, with\enspace\(\mathbb C/{\sim}\cong[0,+\infty)\).

\noindent\textbf{(iii)} On $\mathbb C^*$, define
\[
 z_1\sim z_2\overset{\mathrm{def}}{\iff}
 \arg(z_1)=\arg(z_2).
\]
This is an equivalence relation on $\mathbb C^*$ and
\[
 \mathbb C^*/{\sim}\;\cong\;S^1,
 \qquad S^1=\{z\in\mathbb C\mid |z|=1\}.
\]
% Equality of arguments is understood with a consistent branch/modulo 2pi.

\noindent\textbf{(iv)} On $\mathbb R$, put
\[
 x\sim y\iff x-y\in\mathbb Q.
\]
The manuscript records $\mathbb R/{\sim}\cong\mathbb R$.

\endgroup
; no preamble or page break.
% EDITORIAL NOTES:
% - The opening lines continue the universal property of the quotient set
%   from Page 4. The word attached to the second criterion is overwritten:
%   it may read "inj." or "bij.". The injectivity reading is provisional;
%   if "bij." was intended, surjectivity is also required.
% - The example defining rational numbers writes Z x Z without excluding
%   zero denominators. On all of Z x Z, cross-multiplication is NOT an
%   equivalence relation. The necessary restriction is not silently added.
% - The Cauchy-sequence construction is read as sequences with terms in Q;
%   it is written compactly in the original manuscript.
% - Example (i) includes a small sketch of the quotient projection, here
%   rendered as a commutative diagram.
% - The right-margin scribbles and small shaded circle are not part of an
%   identifiable calculation and are omitted from the main transcription.
% - Example (iv) appears to identify R/~ with R; this is only an abstract
%   (non-canonical) identification, not a specified quotient map.

\subsection*{Seminar 3 - Page 5}
\begingroup
\setlength{\abovedisplayskip}{5pt}
\setlength{\belowdisplayskip}{5pt}
\setlength{\abovedisplayshortskip}{3pt}
\setlength{\belowdisplayshortskip}{3pt}

\noindent\textbf{Universal property (continued).}
For the map $\overline f$ of the preceding page,
\[
\begin{aligned}
 \overline f\text{ is surjective}
   &\iff f\text{ is surjective},\\
 \overline f\text{ is injective}
   &\iff {\sim}\;=R_f.
\end{aligned}
\]
Moreover, for a system of representatives $(x_i)_{i\in I}$,
\[
 M=\bigcup_{i\in I}\widehat{x_i}.
\]

\medskip
\noindent\underline{Examples of quotient sets}

\noindent\textbf{(a) Rational numbers.}
Consider $\mathbb Z\times\mathbb Z$ and the relation
\[
 (m,n)\sim(m',n')
 \overset{\mathrm{def}}{\iff}mn'=nm'.
\]
The notes identify the quotient as
\[
 (\mathbb Z\times\mathbb Z)/{\sim}
   \overset{\mathrm{def}}{=}\mathbb Q,
 \qquad
 r\in\mathbb Q,\quad r=\widehat{(m,n)}.
\]
% Source caveat: n,n' != 0 is required for the intended construction.

\noindent\textbf{(b) Real numbers from Cauchy sequences.}
Let
\[
 \mathcal C=\bigl\{(x_n)_{n\geq1}\text{ with }x_n\in\mathbb Q
       \mid (x_n)_{n\geq1}\text{ is Cauchy}\bigr\}.
\]
Define
\[
 (x_n)_n\sim(y_n)_n
 \overset{\mathrm{def}}{\iff}
 \lim_{n\to\infty}(x_n-y_n)=0.
\]
This is an equivalence relation; with the manuscript's notation,
\[
 \mathcal C/{\sim}\ \overset{\mathrm{not.}}{=}\ \mathbb R,
 \qquad
 r=\widehat{(x_n)_n}\in\mathbb R.
\]

\medskip
\noindent\textbf{(i)} On $\mathbb R$, put
\[
 x\sim y\iff x-y\in\mathbb Z.
\]
This is an equivalence relation and\enspace\(\mathbb R/{\sim}\cong[0,1)\).
The quotient sketch is
\[
\begin{tikzcd}[column sep=large,row sep=small]
 \mathbb R\arrow[r,"p_{\sim}"]\arrow[d,"{x\mapsto x-\lfloor x\rfloor}"']
    & \mathbb R/{\sim} \\
 {[0,1)}\arrow[ur,dashed,"\cong"'] &
\end{tikzcd}
\]
% The floor-map label explicates the diagram's unlabeled vertical arrow.

\noindent\textbf{(ii)} On $\mathbb C$, define
\[
 z_1\sim z_2\overset{\mathrm{def}}{\iff}|z_1|=|z_2|.
\]
This is an equivalence relation on $\mathbb C$, with\enspace\(\mathbb C/{\sim}\cong[0,+\infty)\).

\noindent\textbf{(iii)} On $\mathbb C^*$, define
\[
 z_1\sim z_2\overset{\mathrm{def}}{\iff}
 \arg(z_1)=\arg(z_2).
\]
This is an equivalence relation on $\mathbb C^*$ and
\[
 \mathbb C^*/{\sim}\;\cong\;S^1,
 \qquad S^1=\{z\in\mathbb C\mid |z|=1\}.
\]
% Equality of arguments is understood with a consistent branch/modulo 2pi.

\noindent\textbf{(iv)} On $\mathbb R$, put
\[
 x\sim y\iff x-y\in\mathbb Q.
\]
The manuscript records $\mathbb R/{\sim}\cong\mathbb R$.

\endgroup
; no preamble or page break.
% EDITORIAL NOTES:
% - The opening lines continue the universal property of the quotient set
%   from Page 4. The word attached to the second criterion is overwritten:
%   it may read "inj." or "bij.". The injectivity reading is provisional;
%   if "bij." was intended, surjectivity is also required.
% - The example defining rational numbers writes Z x Z without excluding
%   zero denominators. On all of Z x Z, cross-multiplication is NOT an
%   equivalence relation. The necessary restriction is not silently added.
% - The Cauchy-sequence construction is read as sequences with terms in Q;
%   it is written compactly in the original manuscript.
% - Example (i) includes a small sketch of the quotient projection, here
%   rendered as a commutative diagram.
% - The right-margin scribbles and small shaded circle are not part of an
%   identifiable calculation and are omitted from the main transcription.
% - Example (iv) appears to identify R/~ with R; this is only an abstract
%   (non-canonical) identification, not a specified quotient map.

\subsection*{Seminar 3 - Page 5}
\begingroup
\setlength{\abovedisplayskip}{5pt}
\setlength{\belowdisplayskip}{5pt}
\setlength{\abovedisplayshortskip}{3pt}
\setlength{\belowdisplayshortskip}{3pt}

\noindent\textbf{Universal property (continued).}
For the map $\overline f$ of the preceding page,
\[
\begin{aligned}
 \overline f\text{ is surjective}
   &\iff f\text{ is surjective},\\
 \overline f\text{ is injective}
   &\iff {\sim}\;=R_f.
\end{aligned}
\]
Moreover, for a system of representatives $(x_i)_{i\in I}$,
\[
 M=\bigcup_{i\in I}\widehat{x_i}.
\]

\medskip
\noindent\underline{Examples of quotient sets}

\noindent\textbf{(a) Rational numbers.}
Consider $\mathbb Z\times\mathbb Z$ and the relation
\[
 (m,n)\sim(m',n')
 \overset{\mathrm{def}}{\iff}mn'=nm'.
\]
The notes identify the quotient as
\[
 (\mathbb Z\times\mathbb Z)/{\sim}
   \overset{\mathrm{def}}{=}\mathbb Q,
 \qquad
 r\in\mathbb Q,\quad r=\widehat{(m,n)}.
\]
% Source caveat: n,n' != 0 is required for the intended construction.

\noindent\textbf{(b) Real numbers from Cauchy sequences.}
Let
\[
 \mathcal C=\bigl\{(x_n)_{n\geq1}\text{ with }x_n\in\mathbb Q
       \mid (x_n)_{n\geq1}\text{ is Cauchy}\bigr\}.
\]
Define
\[
 (x_n)_n\sim(y_n)_n
 \overset{\mathrm{def}}{\iff}
 \lim_{n\to\infty}(x_n-y_n)=0.
\]
This is an equivalence relation; with the manuscript's notation,
\[
 \mathcal C/{\sim}\ \overset{\mathrm{not.}}{=}\ \mathbb R,
 \qquad
 r=\widehat{(x_n)_n}\in\mathbb R.
\]

\medskip
\noindent\textbf{(i)} On $\mathbb R$, put
\[
 x\sim y\iff x-y\in\mathbb Z.
\]
This is an equivalence relation and\enspace\(\mathbb R/{\sim}\cong[0,1)\).
The quotient sketch is
\[
\begin{tikzcd}[column sep=large,row sep=small]
 \mathbb R\arrow[r,"p_{\sim}"]\arrow[d,"{x\mapsto x-\lfloor x\rfloor}"']
    & \mathbb R/{\sim} \\
 {[0,1)}\arrow[ur,dashed,"\cong"'] &
\end{tikzcd}
\]
% The floor-map label explicates the diagram's unlabeled vertical arrow.

\noindent\textbf{(ii)} On $\mathbb C$, define
\[
 z_1\sim z_2\overset{\mathrm{def}}{\iff}|z_1|=|z_2|.
\]
This is an equivalence relation on $\mathbb C$, with\enspace\(\mathbb C/{\sim}\cong[0,+\infty)\).

\noindent\textbf{(iii)} On $\mathbb C^*$, define
\[
 z_1\sim z_2\overset{\mathrm{def}}{\iff}
 \arg(z_1)=\arg(z_2).
\]
This is an equivalence relation on $\mathbb C^*$ and
\[
 \mathbb C^*/{\sim}\;\cong\;S^1,
 \qquad S^1=\{z\in\mathbb C\mid |z|=1\}.
\]
% Equality of arguments is understood with a consistent branch/modulo 2pi.

\noindent\textbf{(iv)} On $\mathbb R$, put
\[
 x\sim y\iff x-y\in\mathbb Q.
\]
The manuscript records $\mathbb R/{\sim}\cong\mathbb R$.

\endgroup
